\section*{الفصل \ref{ch:b3:lp} --- فضاءات $L^p$}

\begin{solution}{exo:b3:lp:1}
(a) هي حالة $p = q = 2$ في \cref{thm:b3:lp:holder}: $\abs{\int f\bar
g\,\dd\mu} \leq \int\abs{fg} \leq \norm f_2\norm g_2$ --- أي
كوشي--شوارتز، مع التساوي إذا وفقط إذا كان $\abs f, \abs g$
متناسبين وكانت الأطوار متوافقة.

(b) على فضاء احتمالي، ومن أجل $p \leq q$: نطبّق ينسن
(\cref{exo:b3:lp:10}) بالدالة المحدّبة $\Phi(t) =
\abs t^{q/p}$ على الدالة $\abs f^p$:
$\bigl(\int\abs f^p\bigr)^{q/p} \leq \int\abs f^q$، أي
$\norm f_p \leq \norm f_q$.

(c) يكون $\int\abs{fg} = \norm f_1\norm g_\infty$ إذا وفقط إذا كان $\abs{g} =
\norm g_\infty$ في كل مكان تقريبًا على
$\{f \neq 0\}$ (إذ يجب أن تكون المتراجحة $\abs{fg} \leq \abs f\norm g_\infty$ تساويًا في كل
مكان تقريبًا).
\end{solution}

\begin{solution}{exo:b3:lp:2}
يكون $\int_0^1 x^{-p/2}\dd x < \infty$ إذا وفقط إذا كان $p < 2$: فالأولى في $L^p$ من أجل
$p \in \intco12$. ويكون $\int_1^\infty x^{-p/2}\dd x <
\infty$ إذا وفقط إذا كان $p > 2$:
فالثانية من أجل $p \in \intoo2\infty$ (ومن أجل $p = \infty$: فهي محدودة ---
ونضمّها). والثالثة: من أجل $p <
2$، تُهيمَن بالمقدار $x^{-p/2}$:
فهي \omterm{def:b3:lebesgue:l1}{قابلة للمكاملة}؛ ومن أجل $p = 2$، نعوّض $u = -\ln x$:
$\int_0^1\frac{\dd x}{x(1 + \abs{\ln
x})^2} = \int_0^\infty\frac{\dd u}{(1 + u)^2} < \infty$؛ ومن أجل $p > 2$ تهيمن القوة: فتتباعد. ومنه
$p \in \intcc12$. والرابعة: $\int_\R\frac{\dd x}{(1 + \abs x)^p} < \infty$ إذا وفقط إذا كان $p >
1$؛ وهي
محدودة، ومنه فهي في $L^\infty$ أيضًا: $p \in
\intoc1\infty$. والعبرة: أن
القابلية للمكاملة عند $0$ تحب $p$ الصغير، وعند $\infty$ تحب $p$
الكبير؛ وبضمّ العائقين لا يوجد أي احتواء بين فضاءات $L^p(\R)$.
\end{solution}

\begin{solution}{exo:b3:lp:3}
(a) $\norm{f_n}_p^p = \lambda(I_n) \to 0$ (فعند المستوى الثنائي $k$ يكون الطول $2^{-k}$). لكن
كل $x$ يقع في فترة واحدة من كل مستوى ثنائي: ومنه $f_n(x) = 1$
عددًا لانهائيًّا من المرات و $= 0$ عددًا لانهائيًّا من المرات
(في فترات المستوى نفسه التي لا تحتوي $x$): فلا تقارب في أي
نقطة.

(b) المتتالية $f_{n_k} = \mathbf 1_{\intcc0{2^{-k}}}$ (أي الفترة الأولى من كل مستوى) تتقارب إلى
$0$ عند كل $x > 0$: أي في كل مكان تقريبًا.

(c) في كل مكان تقريبًا دون $L^1$: $n\mathbf 1_{\intoo0{1/n}} \to 0$ في كل مكان تقريبًا،
وتكاملها $1$. وفي $L^1$ دون أي $L^p$ (حيث $p > 1$): $g_n =
\eu^n\,\mathbf 1_{(0,\ \eu^{-n}/n)}$:
فيكون $\norm{g_n}_1 = \frac1n
\to 0$، بينما $\norm{g_n}_p^p = \eu^{(p-1)n}/n \to \infty$ من أجل كل $p > 1$.
\end{solution}

\begin{solution}{exo:b3:lp:4}
الحدّ الأعلى: $\norm f_p \leq \mu(X)^{1/p}\norm f_\infty$
(\cref{prop:b3:lp:inclusions}(a) مع $q = \infty$)، و $\mu(X)^{1/p} \to 1$. والحدّ الأدنى: من
أجل $\varepsilon > 0$، تحقق $A =
\{\abs f > \norm f_\infty - \varepsilon\}$ أن $\mu(A) > 0$ (بتعريف السوپريموم
الجوهري)، و
\[
\norm f_p \geq \Bigl(\int_A \abs f^p\Bigr)^{1/p}
\geq (\norm f_\infty - \varepsilon)\,\mu(A)^{1/p}
\xrightarrow[p\to\infty]{} \norm f_\infty - \varepsilon .
\]
\end{solution}

\begin{solution}{exo:b3:lp:5}
(a) مرتين: فالتحويل $\norm{\tau_hg - g}_p \leq
C^{1/p}\norm{\tau_hg - g}_\infty$ (بحامل منتهي \omterm{def:b3:measure:measure}{القياس}) لا ينحلّ من أجل
$p = \infty$ إلا من حيث إخفاق \emph{كثافة $\mathcal C_c$}
هناك --- وتلك هي الفجوة الحقيقية: إذ ليس للخطوة (2) من
\cref{thm:b3:lp:density} نظيرٌ في $L^\infty$.

(b) إذا كان للدالة $f$ ممثّل متصل بانتظام $g$: فإن $\norm{\tau_hf - f}_\infty = \sup_x\abs{g(x - h) - g(x)} \to
0$.
وبالعكس، لنفترض أن $\norm{\tau_hf - f}_\infty \to 0$. تكون التمليفات $f_\varepsilon = f * \rho_\varepsilon$ متصلة، ويكون
\[
\norm{f_\varepsilon - f}_\infty
\leq \sup_{\norm y \leq \varepsilon}\norm{\tau_yf -
f}_\infty \longrightarrow 0
\]
(إذ الالتفاف متوسطٌ للمنسحبات). وكل $f_\varepsilon$ متصلة
بانتظام ($\norm{\tau_hf_\varepsilon - f_\varepsilon}_\infty \leq
\norm{\tau_hf - f}_\infty$، بالتوسيط)، والنهاية المنتظمة لدوال متصلة
بانتظام هي كذلك: ومنه توافق $f$ في كل مكان تقريبًا دالةً متصلة
بانتظام.
\end{solution}

\begin{solution}{exo:b3:lp:6}
بهولدر، يكون المقدار المكامَل $y \mapsto f(x -
y)g(y)$ من أجل كل $x$ في $L^1$ مع
$\abs{(f*g)(x)} \leq \norm f_p\norm
g_q$: فهي معرَّفة في كل نقطة ومحدودة. وأما الاتصال المنتظم
($p <
\infty$):
\[
\abs{(f*g)(x + h) - (f*g)(x)}
= \Bigl|\int\bigl(\tau_{-h}f - f\bigr)(x - y)\,g(y)\dd y\Bigr|
\leq \norm{\tau_{-h}f - f}_p\,\norm g_q
\xrightarrow[h\to0]{} 0,
\]
بانتظام في $x$ (\cref{thm:b3:lp:density}(3)). وإذا كان $p =
\infty$، فإن $q = 1$:
نكتب $f * g = g * f$ ونجري الحد نفسه مع فعل الانسحاب على
$g \in L^1$.
\end{solution}

\begin{solution}{exo:b3:lp:7}
نثبّت $\chi \in \mathcal C_c^\infty(\intoo ab)$ تحقق $\int\chi =
1$، ونضع $c = \int f\chi$. ولتكن $\varphi \in \mathcal
C_c^\infty$
كيفما اتفق ولتكن $\psi = \varphi -
\bigl(\int\varphi\bigr)\chi$: عندئذٍ $\int\psi = 0$، ومنه تعرّف
$\Phi(x)
= \int_a^x\psi$ دالةً $\Phi \in \mathcal C_c^\infty(\intoo
ab)$ (تنعدم بجوار الطرفين: بجوار $a$ بداهةً،
وبجوار $b$ لأن التكامل الكلي $0$) مع $\Phi' = \psi$. وتعطي
الفرضية أن $\int f\psi = \int f\Phi' = 0$، ومنه
\[
\int f\varphi = \Bigl(\int\varphi\Bigr)\int f\chi = \int
c\,\varphi
\quad\text{من أجل كل} \varphi:
\qquad \int(f - c)\varphi = 0 .
\]
وحسب \cref{cor:b3:lp:fundlemma} (موضَّعةً على $\intoo ab$)، يكون $f =
c$ في كل
مكان تقريبًا.
\end{solution}

\begin{solution}{exo:b3:lp:8}
ليكن $3\delta < d(K, \R^d\setminus U)$ (وهو موجب: \cref{exo:b3:topology:6}(b)؛ وإذا كان $U = \R^d$
صلح أي $\delta$)، ولتكن $K_\delta = \{x : d(x, K) \leq \delta\}$، و
\[
\varphi = \mathbf 1_{K_\delta} * \rho_{\delta/2} .
\]
عندئذٍ $\varphi \in \mathcal C^\infty$ (\cref{thm:b3:lp:regularization}(1)؛ لأن الدالة المميّزة في $L^1$)،
و $0 \leq \varphi \leq 1$ (لأن $\int\rho = 1$)، و $\varphi = 1$ على $K$ (فمن أجل
$x \in K$ يكون $\bar B(x, \delta/2) \subseteq K_\delta$، ومنه يكامل الالتفاف $\rho$ كاملةً)،
و $\operatorname{supp}\varphi \subseteq K_{3\delta/2} \subseteq
U$: وحاملها متراص (لأن $K_\delta$ محدودة). وأما تجزئة
الوحدة: فبإعطاء $K \subseteq U_1\cup\dots\cup U_m$، نختار (بالتراص) متراصات $K_i \subseteq U_i$ تحقق
$K \subseteq
\bigcup \mathring K_i$، ونأخذ $\varphi_i$ كما سبق من أجل $(K_i,
U_i)$، ونضع $\psi_i = \varphi_i\prod_{j <i}(1 -
\varphi_j)$:
فتكون كل $\psi_i \in \mathcal C_c^\infty(U_i)$، ويكون $\sum_i\psi_i = 1 - \prod_i(1 - \varphi_i) = 1$ على $K$.
\end{solution}

\begin{solution}{exo:b3:lp:9}
(a) أُسّ الاستيفاء الموافق للثلاثي $(p_0, q_0) = (1, \infty)$ عند $r = p$ هو
$\alpha = \frac1p$: فتعطي \cref{prop:b3:lp:inclusions}(c) أن $\norm f_p \leq \norm
f_1^{1/p}\norm f_\infty^{1 - 1/p}$.

(b) بأخذ اللوغاريتمات في \cref{prop:b3:lp:inclusions}(c): $\ln\norm f_r \leq
\alpha\ln\norm f_p + (1-\alpha)\ln\norm f_q$ حيث يكون
$\frac1r$ التركيبةَ المحدّبة نفسها للمقدارين $\frac1p, \frac1q$: ومنه
$\frac1p
\mapsto \ln\norm f_p$ محدّبة. ومثال يقع فيه الانتماء إلى $L^p$ على $\intoo{p_0}{p_1}$
بالضبط:
\[
f(x) = x^{-1/p_1}\,\mathbf 1_{\intoo01}(x) +
x^{-1/p_0}\,\mathbf 1_{\intco1\infty}(x):
\]
فالحدّ الأول في $L^p$ إذا وفقط إذا كان $p < p_1$، والثاني إذا
وفقط إذا كان $p >
p_0$.
\end{solution}

\begin{solution}{exo:b3:lp:10}
ليكن $m = \int f\,\dd\mu \in \R$. يوفّر التحدّب مستقيمَ إسناد عند $m$: أي يوجد $s$
يحقق $\Phi(t) \geq \Phi(m) + s(t -
m)$ من أجل كل $t$ (بأخذ $s$ بين المشتقتين من الجانبين،
وهما موجودتان من أجل الدوال المحدّبة). ثم نعوّض $t = f(x)$
ونكامل بالنسبة إلى \omterm{def:b3:measure:measure}{القياس} الاحتمالي $\mu$:
\[
\int\Phi\circ f\,\dd\mu \geq \Phi(m) + s\Bigl(\int f - m\Bigr)
= \Phi\Bigl(\int f\,\dd\mu\Bigr)
\]
(وأما القابلية \omterm{def:b3:measure:measure}{للقياس}: فلأن $\Phi$ متصلة؛ وأما قابلية الجزء
السالب من $\Phi\circ f$ للمكاملة فيضمنها مستقيم الإسناد). وأما
متراجحة الوسطين الحسابي والهندسي: فعلى مجموعة منتهية بأوزان
$w_i$، نأخذ $\Phi = \exp$ و $f = \sum(\ln a_i)\mathbf 1_i$: $\exp\bigl(\sum w_i\ln a_i\bigr) \leq \sum w_ia_i$، أي $\prod a_i^{w_i} \leq \sum w_ia_i$. وأما
رتابة المعيار فهي \cref{exo:b3:lp:1}(b).
\end{solution}

\begin{solution}{exo:b3:lp:11}
(a) نفترض أولًا أن $p, q, r < \infty$ وأن $f, g \geq 0$ (وإلا استعضنا
بالقيم المطلقة). تحقق الأُسس الثلاثة $r$ و $\alpha = \frac{pr}{r-p}$ و $\beta = \frac{qr}{r-q}$ أن
$\frac1r + \frac1\alpha + \frac1\beta = \frac1r +
\frac1p - \frac1r + \frac1q - \frac1r = 1$ (أي علاقة التحجيم). ثم نشطر، من أجل $x$ مثبَّت،
\[
f(y)g(x-y) = \bigl[f(y)^pg(x-y)^q\bigr]^{1/r}\cdot
f(y)^{1 - p/r}\cdot g(x-y)^{1 - q/r},
\]
وتعطي هولدر بالأُسس الثلاثة أن
\[
(f*g)(x) \leq \Bigl(\int f^pg(x-\cdot)^q\Bigr)^{1/r}
\norm f_p^{\,p/\alpha\cdot\alpha/p}\cdots
\]
وبمزيد من الدقة: العامل الثاني هو $\bigl(\int f^{(1-p/r)\alpha}\bigr)^{1/\alpha} =
\norm f_p^{p(1/p - 1/r)}$ لأن $(1 - \frac pr)\alpha = p$،
وبالمثل العامل الثالث هو $\norm g_q^{q(1/q - 1/r)}$. ثم نرفع إلى القوة $r$ ونكامل
في $x$ (بتونيلي على العامل الأول):
\[
\norm{f*g}_r^r \leq \norm f_p^p\,\norm g_q^q\cdot
\norm f_p^{\,rp(1/p - 1/r)}\,\norm g_q^{\,rq(1/q - 1/r)}
= \norm f_p^r\,\norm g_q^r .
\]
وأما الحالات الحدّية ($r = \infty$ أو أن يساوي أحد الأُسس حدّه)
فهي هولدر صرفة أو تقديرات مباشرة.

(b) الحالة $r = \infty$ تفرض $q = p'$: $\abs{f*g(x)} \leq \norm
f_p\norm g_{p'}$ --- أي هولدر بعد
الانسحاب والانعكاس. والحالة $q = 1$ تعطي $r = p$: $\norm{f*g}_p \leq \norm
g_1\norm f_p$، وهي
حصان التمليف (أي محرّك \cref{thm:b3:lp:regularization}). والحالة $p = q = 1$ تعطي
$r = 1$: أي الجبر الالتفافي (\cref{thm:b3:product:convolution}).

(c) نستعيض عن $f, g$ بالدالتين $f_\lambda = f(\lambda\cdot)$ و $g_\lambda = g(\lambda\cdot)$: عندئذٍ
$f_\lambda*g_\lambda =
\lambda^{-d}(f*g)(\lambda\cdot)$، وبمقارنة النظائم،
\[
\text{الطرف الأيسر} \sim \lambda^{-d - d/r}, \qquad
\text{الطرف الأيمن} \sim \lambda^{-d/p - d/q} :
\]
فمتراجحة صحيحة من أجل كل $f, g$ تفرض توافق أُسّي التحجيم، أي
$1 + \frac1r = \frac1p + \frac1q$ بالضبط. وأي تركيبة أخرى تموت عند $\lambda \to 0$ أو
$\infty$: فعلاقة يونغ ليست تسهيلًا بل قانون تحجيم.
\end{solution}

\begin{solution}{exo:b3:lp:12}
(a) نعيِّر $\norm f_p = \norm g_q = 1$. يكامل برهان هولدر متراجحةَ يونغ $\abs{fg} \leq
\frac{\abs f^p}p + \frac{\abs g^q}q$؛
ويفرض تساوي التكاملين تساويًا في كل مكان تقريبًا في يونغ، وهو
يعني (بالتحدّب الأكيد للدالة $\exp$؛ إذ التساوي يقع إذا وفقط
إذا كان $a^p = b^q$) أن $\abs
f^p = \abs g^q$ في كل مكان تقريبًا. وبالرجوع عن
التعيير: $\abs f^p\norm g_q^q = \abs g^q\norm f_p^p$ في كل مكان تقريبًا --- أي التناسب.

(b) مينكوفسكي تطبيقان لهولدر على $\abs{f + g}^{p-1}
\abs f$ و $\abs{f+g}^{p-1}\abs g$؛ ويفرض
التساوي تناسبَ (a) في كليهما: أي إن $\abs f^p$ و $\abs g^p$
يتناسب كلٌّ منهما مع $\abs{f+g}^{(p-1)q} = \abs{f+g}^p$، ومنه $\abs g = t\abs f$ في كل مكان تقريبًا
من أجل ثابت $t \geq 0$؛ كما يجب أن تكون متراجحة المثلث النقطية
الأولى $\abs{f + g} \leq
\abs f + \abs g$ تساويًا في كل مكان تقريبًا أيضًا، وهو يعني في
الحالة العقدية أن للدالتين $f$ و $g$ العمدةَ نفسها في كل مكان
تقريبًا حيث لا تنعدمان. وبالضمّ: $g = tf$ في كل مكان تقريبًا مع
$t > 0$ (إذ كلتاهما غير معدومة).

(c) الحالة $p = 1$: يقع التساوي في $\int\abs{f + g} = \int\abs f +
\int\abs g$ كلما كان للدالتين
$f, g$ نسق الإشارة نفسه (أي العمدة نفسها في كل مكان تقريبًا)
--- دون حاجة إلى أي تناسب: فتصلح $f = \mathbf 1_{\intcc01}$ و $g = \mathbf 1_{\intcc02}$. والحالة
$p = \infty$: يقع $\norm{f+g}_\infty = \norm f_\infty +
\norm g_\infty$ بمجرد أن تبلغ الدالتان ذروتيهما
توافقيًّا عند نقطة مشتركة (أو على امتداد متتالية مشتركة): إذ
يكفي أن تكون $f = g\,$ بجوار نقطة واحدة أيًّا كان سلوكهما في
غيرها. والتحدّب الأكيد لكرات $L^p$ من أجل $1 <
p < \infty$ --- وإخفاقه
عند الطرفين --- هو بالضبط ما تشهد به حالات التساوي هذه.
\end{solution}

\begin{solution}{pb:b3:lp:1}
\textbf{1.} للدالة $f$ حاملٌ في فترة ما $[\alpha, \beta]
\subseteq \intoo0\infty$، ومنه $F = 0$
على $[0, \alpha]$ و $F
\equiv F(\beta)$ على $[\beta, \infty)$: فتنعدم $Hf$ بجوار $0$
وتكون من رتبة $O(1/x)$ عند اللانهاية؛ ويكون $\int_1^\infty x^{-p}\dd x <
\infty$ من أجل
$p > 1$، و $Hf$ متصلة: ومنه $Hf \in L^p$.

\textbf{2.} $\bigl(x^{1-p}F(x)^p\bigr)' = (1-p)x^{-p}F^p +
p\,x^{1-p}F^{p-1}f$. وتنعدم القيمتان الحدّيتان: عند $0$ لأن
$F = 0$ بجوار $0$؛ وعند $\infty$ لأن $x^{1-p}F^p \leq
F(\beta)^p x^{1-p} \to 0$ (حيث $p > 1$).
وبمكاملة المتطابقة على $\intoo0\infty$:
\[
0 = (1 - p)\int_0^\infty\Bigl(\frac Fx\Bigr)^p\dd x
+ p\int_0^\infty\Bigl(\frac Fx\Bigr)^{p-1}f(x)\,\dd x,
\]
وهي العلاقة المعروضة.

\textbf{3.} هولدر بالأُسّين $q = \frac p{p-1}$ و $p$:
\[
\int\Bigl(\frac Fx\Bigr)^{p-1}f
\leq \Bigl(\int\Bigl(\frac
Fx\Bigr)^{p}\Bigr)^{1 - 1/p}\,\norm f_p ,
\]
ومنه $\norm{Hf}_p^p \leq \frac p{p-1}\norm{Hf}_p^{p-1}\norm
f_p$؛ ثم نقسم (والمقدار منتهٍ حسب السؤال 1، وإن كان $0$
فلا شيء نبرهن عليه).

\textbf{4.} ليكن $f \in L^p$ مع $f \geq 0$. نختار $\varphi_k \in
\mathcal C_c(\intoo0\infty)$ تحقق
$\varphi_k \to f$ في $L^p$ (\cref{thm:b3:lp:density}(2)، مقاطعةً مع نصف المستقيم
المفتوح --- بتقريب $f\mathbf 1_{[1/k, k]}$ ثم بالتقطير)، ونستعيض عن $\varphi_k$
بالمقدار $\abs{\varphi_k}$ (وهو لا يزال \omterm{def:b3:topology:continuity}{متصلًا}، وأقرب إلى $f \geq 0$).
ومن أجل كل $x > 0$ مثبَّت، تعطي هولدر على $\intoo0x$ أن
\[
\abs{H\varphi_k(x) - Hf(x)} \leq
\frac1x\,x^{1 - 1/p}\,\norm{\varphi_k - f}_p \to 0 :
\]
أي $H\varphi_k \to Hf$ نقطةً نقطة. وبفاتو وبالسؤال 3:
\[
\int (Hf)^p \leq \liminf_k\int(H\varphi_k)^p
\leq \Bigl(\frac p{p-1}\Bigr)^p\liminf_k\norm{\varphi_k}_p^p
= \Bigl(\frac p{p-1}\Bigr)^p\norm f_p^p .
\]
وأما من أجل $f$ ذات إشارة أو عقدية: فإن $\abs{Hf} \leq H\abs f$ نقطةً نقطة،
وتنطبق الحالة غير السالبة على $\abs f$.

\textbf{5.} $\norm{f_A}_p^p = \int_1^A t^{-1}\dd t = \ln A$. ومن أجل $1 \leq x \leq A$:
\[
(Hf_A)(x) = \frac1x\int_1^x t^{-1/p}\dd t
= \frac{x^{1 - 1/p} - 1}{(1 - \tfrac1p)\,x}
= \frac p{p-1}\,x^{-1/p}\bigl(1 - x^{-(1 - 1/p)}\bigr).
\]

\textbf{6.} نثبّت $\varepsilon > 0$ و $X_0$ بحيث يكون $(1 -
x^{-(1-1/p)})^p \geq 1 - \varepsilon$ من أجل
$x \geq X_0$. عندئذٍ
\[
\norm{Hf_A}_p^p \geq \int_{X_0}^A\Bigl(\frac
p{p-1}\Bigr)^p\frac{1 - \varepsilon}{x}\,\dd x
= \Bigl(\frac p{p-1}\Bigr)^p(1 - \varepsilon)\,(\ln A - \ln
X_0),
\]
ومنه $\dfrac{\norm{Hf_A}_p^p}{\norm{f_A}_p^p} \geq \bigl(\frac
p{p-1}\bigr)^p(1 - \varepsilon)\bigl(1 - \frac{\ln X_0}{\ln
A}\bigr) \to \bigl(\frac p{p-1}\bigr)^p(1 - \varepsilon)$ حين $A \to \infty$: فالثابت غير قابل للتحسين.

\textbf{7.} يفرض التساوي في السؤال 3 تساويًا في هولدر: أي أن
يتناسب $f^p$ مع $\bigl(\frac Fx\bigr)^{(p-1)q}
= \bigl(\frac Fx\bigr)^p$ في كل مكان تقريبًا، أي $f = \gamma\,\frac Fx$ في كل
مكان تقريبًا من أجل $\gamma \geq 0$ ما. وبما أن $F(x) = \int_0^xf$ متصلة
اتصالًا مطلقًا مع $F' = f$ في كل مكان تقريبًا، فإن $F$ تحل
$F' =
\gamma F/x$: فعلى أي فترة تكون فيها $F > 0$، يكون $(\ln F)' =
\gamma/x$، ومنه
$F = c\,x^{\gamma}$ و $f = c\gamma
x^{\gamma - 1}$ هناك. لكن لا تنتمي أي قوة غير معدومة
$x^{\gamma-1}$ إلى $L^p(\intoo0\infty)$ (إذ يلزم $p(\gamma - 1) < -1$ عند $\infty$
و $> -1$ عند $0$: وهما متنافيان)، ولا يمكن أن تنعدم $F$ انعدامًا
تامًّا إلا إذا كان $f = 0$. ومنه يقتضي التساوي أن $f =
0$.

\textbf{8.} بإعطاء $(a_n)$ متناقصة بالمعنى الواسع وغير سالبة $\geq 0$،
نعرّف الدالة السلّمية $g(t) = a_{\lceil t\rceil}$ على $\intoo0{+\infty}$: فهي متناقصة بالمعنى
الواسع، مع $\int_0^\infty g^p =
\sum_na_n^p$ و $\int_0^ng = a_1 + \dots + a_n$، ومنه $(Hg)(n) = \frac{a_1 + \dots + a_n}n$. ومتوسط دالة متناقصة
بالمعنى الواسع متناقص بالمعنى الواسع، ومنه كذلك $Hg$، و
\[
\sum_{n\geq1}\Bigl(\frac{a_1{+}\dots{+}a_n}n\Bigr)^p
= \sum_{n\geq1}(Hg)(n)^p
\leq \sum_{n\geq1}\int_{n-1}^n (Hg)(t)^p\,\dd t
= \norm{Hg}_p^p
\leq \Bigl(\frac p{p-1}\Bigr)^p\sum_n a_n^p
\]
حسب الجزء الأول. وأما من أجل متتالية غير سالبة عامة، فلتكن
$(a_n^*)$ إعادةَ ترتيبها المتناقصة بالمعنى الواسع (وهي ممكنة
حين $a_n \to 0$، ولنا أن نفترض ذلك --- وإلا فالطرفان لانهائيان):
فلا يتغيّر الطرف الأيمن، ويكون كل مجموع جزئي $a_1 + \dots
+ a_n$ أصغر من
أو يساوي $a_1^* + \dots + a_n^*$ (أي الحدود الكبرى $n$): ومنه لا يزيد الطرف
الأيسر إلا نموًّا. فتصح المتراجحة إذًا من أجل كل $(a_n)$.

\textbf{9.} إذا كان $(a_n) \in \ell^p$، فإن متتالية متوسطات تشيزارو في
$\ell^p$ بمعيار $\leq \frac p{p-1}\norm
a_p$. ومثال (حيث $p = 2$): $a_n = \frac1{\sqrt n\,\ln n}$ (أي
$n
\geq 2$): فيكون $\sum a_n^2 = \sum\frac1{n\ln^2n} < \infty$، ومع ذلك $\sum a_n = \infty$ (باختبار التكامل)؛
ومع ذلك تظل هاردي تضمن $\sum_n\bigl(\frac{a_1 + \dots + a_n}n\bigr)^2 < \infty$.

\textbf{10.} من أجل $f = \mathbf 1_{\intcc01}$: يكون $Hf(x) = 1$ على $\intoc01$
و $= \frac1x$ من أجل $x \geq 1$: ومنه $\int_0^\infty Hf =
1 + \int_1^\infty\frac{\dd x}x = \infty$، بينما $\norm f_1 =
1$.
وينهار البرهان عند نقطتين: إذ ينفجر الثابت $\frac p{p-1}$ حين
$p \to 1$، ولم يعد الحدّ الحدّي $x^{1-p}F^p$ ينعدم عند اللانهاية
من أجل $p = 1$. فمتراجحة هاردي ظاهرة صادقة من ظواهر $p > 1$.

\textbf{11.} نختار أطول فترة $B_{i_1}$؛ ونطرح كل فترة تقاطعها؛
ثم نختار أطول ناجٍ $B_{i_2}$؛ ونكرّر (والفترات منتهية العدد).
فتكون المختارة منفصلة بحكم البناء، وكل فترة $B$ مطروحة قد قاطعت
فترةً مختارةً لا تقلّ عنها طولًا: وفترةٌ تقاطع فترةً أطول منها
أو مساويةً لها محتواةٌ في ثلاثة أمثالها، $B \subseteq
3B_{i_l}$.

\textbf{12.} المجموعة $\{Mf > t\}$ مفتوحة: إذ كل متوسط $x \mapsto
\frac1{2r}\int_{x-r}^{x+r}\abs f$
متصل (بالتقارب المهيمن في $x$)، وسوپريموم دوال متصلة نصفُ متصل
من الأسفل. وكل $x$ فيها يملك $B_x =
\intoo{x-r_x}{x+r_x}$ يحقق $\int_{B_x}\abs f >
t\,\lambda(B_x)$. ومن أجل
متراصة $K \subseteq \{Mf > t\}$: يغطّي $K$ عددٌ منتهٍ من الفترات $B_x$، وتستخرج
فيتالي (السؤال 11) فترات منفصلة $B_1', \dots, B_k'$ تحقق $K \subseteq
\bigcup_l3B_l'$، ومنه
\[
\lambda(K) \leq 3\sum_l\lambda(B_l') <
\frac3t\sum_l\int_{B_l'}\abs f \leq \frac3t\,\norm f_1
\]
بالانفصال؛ ويخلص الانتظام الداخلي (\cref{thm:b3:measure:regularity}) إلى النتيجة.

\textbf{13.} من أجل $x > 1$: مع $r \in \intcc{x-1}x$، يكون المتوسط $\frac{r - x + 1}{2r}$،
وهو متزايد في $r$؛ ومن أجل $r
\geq x$ يكون $\frac1{2r}$، وهو
متناقص: فالسوپريموم هو $\frac1{2x}$، وهو مبلوغ عند $r = x$.
ومنه $M\mathbf 1_{\intcc01}
\notin L^1$. وعمومًا، إذا كان $\int_I\abs f = c > 0$ على فترة محدودة $I \subseteq \intcc{-C}C$،
فإن $Mf(x) \geq
\frac{c}{2(\abs x + C)}$ من أجل كل $x$: فلا تكون \omterm{def:b3:lebesgue:l1}{قابلة للمكاملة} أبدًا إلا إذا
كان $f = 0$ في كل مكان تقريبًا.

\textbf{14.} مع $f_t = f\,\mathbf 1_{\abs f > t/2}$: $M(f - f_t) \leq \frac t2$، ومنه $\{Mf > t\} \subseteq \{Mf_t
> \frac t2\}$، ويعطي السؤال 12
أن $\lambda(Mf > t) \leq
\frac6t\int_{\abs f > t/2}\abs f$. وبكعكة الطبقات (\cref{prop:b3:product:layercake}) وتونيلي:
\[
\norm{Mf}_p^p = p\int_0^\infty t^{p-1}\lambda(Mf > t)\dd t
\leq 6p\int\abs f\int_0^{2\abs f}t^{p-2}\,\dd t\,\dd\lambda
= \frac{6p\,2^{p-1}}{p-1}\,\norm f_p^p .
\]

\textbf{15.} نكتب $A_rf(x) = \frac1{2r}\int_{x-r}^{x+r}f$. ومن أجل $g$ متصلة: $A_rg(x) \to g(x)$ في كل
مكان. وبإعطاء $\varepsilon > 0$، نشطر $f = g + h$ حيث $g$ متصلة ذات حامل
متراص و $\norm h_1 < \varepsilon$ (\cref{thm:b3:lp:density})؛ عندئذٍ
\[
\limsup_{r\to0}\,\abs{A_rf(x) - f(x)} \leq Mh(x) +
\abs{h(x)},
\]
ومنه يكون \omterm{def:b3:measure:measure}{قياس} $\Omega_\delta = \{\limsup_r\abs{A_rf - f} > \delta\}
\subseteq \{Mh > \tfrac\delta2\} \cup \{\abs h >
\tfrac\delta2\}$ أصغر من أو يساوي $\leq \frac{6\varepsilon}\delta
+ \frac{2\varepsilon}\delta$ (بالسؤال 12؛
وبماركوف). و $\varepsilon$ كيفما اتفق: $\lambda(\Omega_\delta) = 0$؛ وبالاتحاد على
$\delta = \frac1k$: يكون $A_rf \to f$ في كل مكان تقريبًا.

\textbf{16.} (a) من أجل كل $q \in \Q$، يعطي السؤال 15 مطبَّقًا
على $\abs{f - q}$ أن $A_r\abs{f - q}(x) \to \abs{f(x) - q}$ في كل مكان تقريبًا؛ وعلى تقاطع هذه
المجموعات التامة \omterm{def:b3:measure:measure}{القياس}، نختار $q$ يحقق $\abs{f(x) - q} < \eta$: فيكون $\limsup_rA_r\abs{f - f(x)}(x) \leq 2\eta$
من أجل كل $\eta$: أي إن كل نقطة تقريبًا نقطةُ لوبيغ. (b) وعند
نقطة لوبيغ،
\[
\Bigl|\frac{F(x + h) - F(x)}h - f(x)\Bigr|
= \Bigl|\frac1h\int_x^{x+h}\bigl(f - f(x)\bigr)\Bigr|
\leq 2\,A_{\abs h}\abs{f - f(x)}(x) \to 0 :
\]
أي $F' = f$ في كل مكان تقريبًا --- فالدوال الأصلية لدوال $L^1$
تشتقّ فعلًا رجوعًا؛ والسلّم (\cref{pb:b3:measure:1}) مثال مضاد للاتجاه
\emph{العكسي} وحده.

\textbf{17.} نطبّق السؤال 15 على $\mathbf 1_{A\cap[-n,n]}$ ونترك $n$ ينمو: فمن
أجل $x \in A$ في كل مكان تقريبًا تكون الكثافة $\frac{\lambda(A\cap\intcc{x-r}{x+r})}{2r} \to 1$. وأما
شتاينهاوس: فحول نقطة كثافة نأخذ $r$ بكثافة $> \frac34$؛ ومن أجل
$\abs t < \frac r2$، تملأ كلٌّ من $A$ و $A + t$ أكثر من $\frac32r$ من فترة
طولها $\leq \frac52r$، ومنه تتقاطعان: $\intoo{-r/2}{r/2} \subseteq A - A$.

\textbf{18.} لتكن $G(x) = x^{\alpha+1-p}F(x)^p$: تنعدم $G$ عند $0$ (لأن $F$ تنعدم
بجوار $0$) وعند $\infty$ (لأن $F$ محدودة، و $\alpha + 1 - p < 0$)، ومنه
$\int_0^\infty G' = 0$ مع
\[
G' = (\alpha + 1 - p)\,x^{\alpha-p}F^p +
p\,x^{\alpha+1-p}F^{p-1}f
= (\alpha + 1 - p)\Bigl(\frac Fx\Bigr)^px^\alpha +
p\Bigl(\frac Fx\Bigr)^{p-1}f\,x^\alpha .
\]
ومنه $\int(F/x)^px^\alpha = \frac{p}{p-1-\alpha}
\int(F/x)^{p-1}f\,x^\alpha$؛ وتنهي هولدر بالنسبة إلى \omterm{def:b3:measure:measure}{القياس} $x^\alpha\dd x$
(بالأُسّين $\frac p{p-1}$ و $p$) البرهانَ كما في السؤال 3. وأما
الحالة الحدّية $\alpha = p - 1$: فمع $f(t) =
\frac1t\mathbf 1_{\intcc1A}$، يكون الطرف الأيمن
$\ln A$ بينما يحتوي الأيسر على $\int_1^A\frac{(\ln x)^p}x\dd x =
\frac{(\ln A)^{p+1}}{p+1}$: فلا يصمد أي ثابت حين
$A \to
\infty$.

\textbf{19.} بتونيلي على $\{0 < t < x\}$:
\[
\langle Hf, g\rangle =
\int_0^\infty\frac{g(x)}x\int_0^xf(t)\,\dd t\,\dd x
= \int_0^\infty f(t)\int_t^\infty\frac{g(x)}x\,\dd x\,\dd t
= \langle f, H^*g\rangle .
\]
وأما بالثنوية: $\norm{H^*f}_p = \sup\{\langle f, Hg\rangle : g
\geq 0, \norm g_q \leq 1\} \leq \norm f_p\cdot
\sup\norm{Hg}_q \leq \frac q{q-1}\norm f_p = p\,\norm f_p$، إذ تُستدعى هاردي في $L^q$، وثابتها
$\frac q{q-1}$ يساوي $p$.

\textbf{20.} نشطر على امتداد القطر (المعدوم \omterm{def:b3:measure:measure}{القياس}). فعلى
$\{y \leq
x\}$:
\[
\iint_{y\leq x}\frac{f(x)g(y)}{x}\,\dd y\,\dd x
= \int f(x)\,(Hg)(x)\,\dd x \leq \norm f_p\,\norm{Hg}_q
\leq p\,\norm f_p\norm g_q,
\]
لأن هاردي في $L^q$ تحمل الثابت $\frac q{q-1} =
p$. وبالتناظر، $\iint_{y>x} = \int g\,(Hf) \leq
\norm g_q\norm{Hf}_p \leq q\,\norm f_p\norm g_q$
($\frac
p{p-1} = q$). وفي المجموع: $(p + q)\,\norm f_p\norm g_q$.

\textbf{21.} من أجل $a_n = n^{-1/p}$ حيث $n \leq N$: يكون
الطرف الأيمن $\bigl(\frac p{p-1}\bigr)^p\sum_{n\leq N}\frac1n =
\bigl(\frac p{p-1}\bigr)^p\ln N + O(1)$. وعلى اليسار، من أجل $n
\leq N$: $a_1 + \dots + a_n \geq \int_1^{n+1}t^{-1/p}\dd t =
\frac{p}{p-1}\bigl((n+1)^{1-1/p} - 1\bigr)$،
ومنه يكون متوسط تشيزارو ذو الرتبة $n$ مساويًا $\geq \frac p{p-1}n^{-1/p}(1 -
o(1))$ بانتظام
من أجل $n$ في أي مجال $n \geq n_0(\eta)$؛ وبالرفع إلى القوة $p$ والجمع،
يكون الطرف الأيسر $\geq
\bigl(\frac p{p-1}\bigr)^p(1 - \eta)\ln N + O_\eta(1)$. وبالقسمة وبجعل $N \to \infty$ ثم
$\eta \to 0$: لا يصلح أي ثابت أصغر من $\bigl(\frac p{p-1}\bigr)^p$.

\textbf{22.} كون $H$ محدودًا على $L^p$: أي إن المتوسطات
التراكمية لا تضخّم النظائم من النمط $p$ (بثابت $\frac p{p-1}$)؛
وتشيزارو على $\ell^p$: الشيء نفسه، متقطّعًا؛ وكون $M$ محدودًا
على $L^p$: أي إن أفضل متوسط محلي يظل تحت السيطرة (بثابت
$O(\frac1{p-1})$). وتخفق الثلاثة جميعًا عند $p = 1$، ولسبب
واحد: أن توسيط وحدة كتلة مركَّزة ينتج ذيلًا من الشكل $\frac1x$
(السؤالان 10 و13)، والمقدار $\frac1x$ ينتمي إلى كل $L^p$ بجوار
اللانهاية عدا $L^1$. فالتنعيم ينشر الكتلة إلى حدود القابلية
للمكاملة التوافقية بالضبط.

\textbf{23.} نضع $b_n = a_n^{1/p}$، ومنه $\sum b_n^p = \sum
a_n$. ويعطي السؤال 8
\[
\sum_{n\geq1}\Bigl(\frac{b_1 + \dots + b_n}n\Bigr)^{p}
\leq \Bigl(\frac p{p-1}\Bigr)^{p}\sum_{n\geq1}a_n,
\]
وتحدّ متراجحة الوسطين الحسابي والهندسي كل حدّ من الأسفل:
\[
\Bigl(\frac{b_1 + \dots + b_n}n\Bigr)^{p}
\geq \bigl(b_1\cdots b_n\bigr)^{p/n}
= \bigl(a_1\cdots a_n\bigr)^{1/n}.
\]
ومنه $\sum(a_1\cdots a_n)^{1/n} \leq \bigl(\frac
p{p-1}\bigr)^p\sum a_n$ من أجل \emph{كل} $p > 1$. ومع $m =
p - 1$،
\[
\Bigl(\frac p{p-1}\Bigr)^{p} =
\Bigl(1 + \frac1m\Bigr)^{m+1},
\]
وهو يتناقص إلى $\eu$ حين $m \to \infty$ (وهي المتتالية الرتيبة
الكلاسيكية العليا للعدد $\eu$). وبأخذ الإنفيموم على $p$ نحصل
على متراجحة كارلمان بالثابت $\eu$.

\textbf{24.} من أجل $a_n = \frac1n$، $(a_1\cdots a_n)^{1/n} =
(n!)^{-1/n}$. ويعطي حصر
\cref{pb:b3:product:1} أن $\sqrt{2\pi n}\,(n/\eu)^n \leq n! \leq \sqrt{2\pi
n}\,(n/\eu)^n\eu^{1/(12n)}$، ومنه
\[
(n!)^{1/n} = \frac n\eu\,(2\pi n)^{1/(2n)}
\eu^{O(1/n^2)}
= \frac n\eu\Bigl(1 +
O\Bigl(\frac{\ln n}{n}\Bigr)\Bigr),
\]
لأن $(2\pi n)^{1/(2n)} = \exp\bigl(\frac{\ln(2\pi
n)}{2n}\bigr)$. وبالقلب، $(n!)^{-1/n} = \frac\eu n(1 +
O(\frac{\ln n}n))$، وبالجمع على $n \leq N$:
\[
\sum_{n\leq N}(n!)^{-1/n}
= \eu\ln N + O(1),
\qquad
\eu\sum_{n\leq N}\frac1n = \eu\ln N + O(1)
\]
(إذ تتقارب متسلسلة الخطأ $\sum\frac{\ln n}{n^2}$). ومتراجحة كارلمان بثابت $c$
تفرض $\eu\ln N
+ O(1) \leq c\,(\ln N + O(1))$، ومنه $c \geq \eu$ بالقسمة على $\ln N$. وتصطفّ
المتتاليات المُمثِّلة: فثابت هاردي يُبلَغ بالمتتالية $n^{-1/p}$
(السؤال 21)، وثابت كارلمان بالمتتالية $n^{-1}$ --- وفي كل حالة
تكون المتتالية من النمط التوافقي جالسةً خارج الفضاء المتوسَّط
عليه بقليل.

\textbf{25.} من أجل $f = \mathbf 1_{\intcc01}$: يكون $Hf(x) = 1$ على $\intoc01$
و $Hf(x) = \frac1x$ من أجل $x > 1$، ومنه $\norm{Hf}_2^2 = 1 + \int_1^\infty x^{-2}\dd x = 2$، وتكون النسبة $\sqrt2 \approx 1.414$،
أي نحو $71\%$ من الحدّ الأمثل. وأما من أجل $f_A$ (حيث $p = 2$):
فإن $F(x) = 2(\sqrt x - 1)$ على $\intcc1A$، ومنه على هذا المجال $Hf_A = 2x^{-1/2} - 2x^{-1}$، ومن
أجل $x > A$ يكون $Hf_A(x) = 2(\sqrt A - 1)/x$. وبالتربيع والمكاملة،
\begin{align*}
\int_1^A\bigl(2x^{-1/2} - 2x^{-1}\bigr)^2\dd x
&= 4\ln A - 12 + \frac{16}{\sqrt A} - \frac4A,\\
\int_A^{\infty}\frac{4(\sqrt A - 1)^2}{x^2}\,\dd x
&= 4 - \frac{8}{\sqrt A} + \frac4A,
\end{align*}
ومنه $\norm{Hf_A}_2^2 = 4\ln A - 8 + 8A^{-1/2}$؛ وبالقسمة على $\norm{f_A}_2^2 = \ln A$ نحصل على المتطابقة
المذكورة. وعند $A
= \eu^{10}$: $4 - \frac{8(1 - \eu^{-5})}{10} = 3.2054$، ومنه تكون النسبة $\sqrt{3.2054} \approx 1.790 < 2$.
والنقص $4 - \norm{Hf_A}_2^2/\norm{f_A}_2^2 \sim 8/\ln A$ لا يتناقص إلا لوغاريتميًّا: فلبلوغ النسبة $1.99$
يلزم $\ln A \approx 200$، أي $A \approx 10^{87}$. فالثوابت المثلى مبرهنات لا تجارب.
\end{solution}
